\chapter*{TÓM TẮT}
\addcontentsline{toc}{chapter}{\textbf{TÓM TẮT}}

Sự phát triển mạnh mẽ của Trí tuệ nhân tạo, đặc biệt là các Mô hình ngôn ngữ lớn (Large Language Models --- LLMs), đã mở ra nhiều tiềm năng ứng dụng đột phá trong lĩnh vực giáo dục. Dẫu vậy, việc tích hợp trực tiếp LLM vào giáo dục phổ thông --- tiêu biểu là môn Khoa học tự nhiên cấp Trung học cơ sở, môn học tích hợp ba phân môn Vật lí, Hoá học và Sinh học --- vẫn vấp phải hai rào cản căn bản: hiện tượng sinh thông tin sai lệch (ảo giác --- hallucination) và sự thiếu hụt cơ chế truy xuất minh bạch từ nguồn sách giáo khoa chuẩn mực trong nước.

Đồ án tốt nghiệp này nghiên cứu và xây dựng một hệ thống trợ lý ảo ứng dụng kỹ thuật sinh văn bản tăng cường truy xuất đa phương thức (Multi-modal Retrieval-Augmented Generation --- RAG \cite{lewis2020retrieval}) dành riêng cho môn Khoa học tự nhiên cấp Trung học cơ sở. Kho tri thức của hệ thống được số hoá toàn diện từ \textbf{12 cuốn sách giáo khoa} lớp 6, 7, 8 và 9 thuộc ba bộ sách hiện hành (Kết nối tri thức với cuộc sống, Chân trời sáng tạo, Cánh Diều): \textbf{2.387/2.399 trang} nội dung được lập chỉ mục thành \textbf{16.515 đoạn văn bản} (sử dụng mô hình nhúng \texttt{BAAI/bge-m3} \cite{chen2024bgem3}, $1024$ chiều) và \textbf{3.881 khung hình} (sử dụng \texttt{clip-vit-base-patch16} \cite{radford2021clip}, $512$ chiều) trên hệ quản trị cơ sở dữ liệu vector ChromaDB, song song với một chỉ mục thưa BM25 \cite{robertson2009bm25} đồng bộ định danh đoạn. Mô hình sinh ngôn ngữ được lựa chọn là Qwen2.5-3B-Instruct \cite{qwen25_2024}. Toàn bộ quy trình trích xuất và lập chỉ mục dữ liệu được thiết kế tối ưu để vận hành hiệu quả trực tiếp trên \textbf{CPU} của máy trạm thông thường mà không phụ thuộc vào phần cứng GPU đắt đỏ.

Đóng góp cốt lõi về mặt phương pháp luận của đồ án tập trung vào \textbf{pha truy xuất}. Hệ thống kết hợp chặt chẽ giữa cơ chế định tuyến theo ý định truy vấn, truy xuất lai (kết hợp BM25 và tìm kiếm ngữ nghĩa) cùng bước tái xếp hạng bằng mô hình cross-encoder \texttt{BAAI/bge-reranker-v2-m3}, với các công tắc độc lập cho phép đánh giá chính xác đóng góp của từng thành phần. Bộ dữ liệu kiểm thử được thiết lập thông qua phương pháp lấy mẫu \textbf{ngẫu nhiên trên toàn bộ tập ngữ liệu}, bao gồm \textbf{240 câu hỏi}: $158$ câu hỏi văn bản, $52$ câu hỏi hình ảnh và $30$ câu hỏi \textbf{ngoài phạm vi} môn học nhằm kiểm chứng trực tiếp khả năng nhận diện ngoại vi và từ chối trả lời, tránh hiện tượng phỏng đoán sai lệch. Ở cấu hình đề xuất (truy xuất lai kết hợp tái xếp hạng), hệ thống đạt chỉ số \textbf{MRR $0{,}7789$} trên $210$ câu hỏi có nhãn nguồn, vượt trội so với các cấu hình đối chứng (BM25 thuần đạt MRR $0{,}6636$; ngữ nghĩa thuần đạt MRR $0{,}5664$; phương pháp lai chưa qua tái xếp hạng đạt MRR $0{,}6947$) trên hầu hết các chỉ số, ngoại trừ một hiện tượng nghịch đảo đáng lưu ý: chỉ số $R@10$ của cấu hình đề xuất ($0{,}2458$) thấp hơn so với cấu hình phương pháp lai chưa tái xếp hạng ($0{,}2643$). Ngoài ra, kết quả thực nghiệm cũng xác nhận rằng cơ chế lọc theo khoảng cách tương đối làm suy giảm hiệu năng của truy xuất lai ở quy mô triển khai thực tế (MRR giảm từ $0{,}7779$ xuống $0{,}7157$), trong khi bước tái xếp hạng luôn mang lại cải thiện tích cực trên cả ba kênh truy xuất.

Chất lượng câu trả lời sinh ra được đánh giá độc lập theo phương pháp LLM-as-a-judge \cite{zheng2023llmjudge} thông qua một tập hợp bốn mô hình ngôn ngữ xoay vòng, phân tách chi tiết theo từng loại câu hỏi: nhóm câu hỏi ngoài phạm vi đạt kết quả gần như tuyệt đối (Tính đúng $4{,}87/5$, Độ trung thực và Độ liên quan $5{,}00/5$); nhóm văn bản đạt chất lượng tốt (Tính đúng $4{,}19/5$, Độ trung thực $4{,}56/5$, Độ liên quan $4{,}59/5$); nhóm câu hỏi dựa trên hình ảnh đạt mức thấp hơn (Tính đúng $3{,}08/5$). Hệ thống đảm bảo tính minh bạch học thuật khi thông tin trích dẫn trang sách được thiết lập mang tính \textit{xác định} từ siêu dữ liệu của đoạn văn bản truy xuất, triệt tiêu hoàn toàn khả năng mô hình tự sinh sai lệch số trang. Báo cáo cũng thẳng thắn phân tích các giới hạn kỹ thuật: giải pháp kiến trúc lai giữa Tesseract và mô hình thị giác--ngôn ngữ dành cho công thức Vật lí, Hoá học đã hoàn thành xử lý trên toàn bộ ngữ liệu ($4\,293$ lượt ghép dòng thành công), song mức độ cải thiện của chỉ số dưới vẫn cần được tiếp tục lượng hoá qua các đợt đo đối sánh chuyên sâu; đồng thời thực nghiệm đối chiếu đa phương thức trên $52$ câu hỏi hình ảnh cho thấy ngữ cảnh đa phương thức chưa mang lại cải thiện rõ rệt (Tính đúng $\Delta{-0{,}077}$, Độ trung thực và Độ liên quan $\Delta{-0{,}038}$).

% \vspace{1em}
% \noindent\textbf{Từ khóa:} Trí tuệ nhân tạo, Mô hình ngôn ngữ lớn (LLM), Sinh văn bản tăng cường truy xuất (RAG), RAG đa phương thức, Truy xuất lai, Trợ lý ảo, Khoa học tự nhiên trung học cơ sở, Qwen2.5, Vector Database.
